\documentclass[11pt,a4paper]{article}

% ============================================
% PACKAGES
% ============================================
\usepackage[utf8]{inputenc}
\usepackage[T1]{fontenc}
\usepackage[english]{babel}
\usepackage{geometry}
\usepackage{hyperref}
\usepackage{graphicx}
\usepackage{xcolor}
\usepackage{listings}
\usepackage{tcolorbox}
\usepackage{fontawesome5}
\usepackage{enumitem}
\usepackage{booktabs}
\usepackage{fancyhdr}
\usepackage{titlesec}
\usepackage{amsmath}
\usepackage{array}

% ============================================
% PAGE LAYOUT
% ============================================
\geometry{
    left=2.5cm,
    right=2.5cm,
    top=3cm,
    bottom=3cm,
    headheight=15pt
}

% ============================================
% COLORS
% ============================================
\definecolor{primaryblue}{RGB}{0,102,204}
\definecolor{warningred}{RGB}{220,50,50}
\definecolor{successgreen}{RGB}{40,167,69}
\definecolor{infobox}{RGB}{240,248,255}
\definecolor{warningbox}{RGB}{255,243,205}
\definecolor{codebg}{RGB}{248,249,250}
\definecolor{codeframe}{RGB}{222,226,230}

% ============================================
% HYPERLINKS
% ============================================
\hypersetup{
    colorlinks=true,
    linkcolor=primaryblue,
    filecolor=primaryblue,
    urlcolor=primaryblue,
    citecolor=primaryblue,
    pdftitle={Payroll Analytics System - User Manual},
    pdfauthor={MARS Project},
    pdfsubject={Installation and User Guide},
    pdfkeywords={Payroll, Analytics, Forecasting, Dashboard}
}

% ============================================
% HEADERS AND FOOTERS
% ============================================
\pagestyle{fancy}
\fancyhf{}
\fancyhead[L]{\small\textit{Payroll Analytics System}}
\fancyhead[R]{\small\textit{User Manual}}
\fancyfoot[C]{\thepage}
\renewcommand{\headrulewidth}{0.4pt}
\renewcommand{\footrulewidth}{0.4pt}

% ============================================
% SECTION FORMATTING
% ============================================
\titleformat{\section}
{\color{primaryblue}\Large\bfseries}
{\thesection}{1em}{}

\titleformat{\subsection}
{\color{primaryblue}\large\bfseries}
{\thesubsection}{1em}{}

\titleformat{\subsubsection}
{\color{primaryblue}\normalsize\bfseries}
{\thesubsubsection}{1em}{}

% ============================================
% CUSTOM BOXES
% ============================================
\tcbuselibrary{skins,breakable}

\newtcolorbox{infobox}[1][]{
    colback=infobox,
    colframe=primaryblue,
    fonttitle=\bfseries,
    title=\faInfoCircle\ Information,
    breakable,
    enhanced,
    attach boxed title to top left={yshift=-2mm, xshift=5mm},
    boxed title style={colback=primaryblue, colframe=primaryblue},
    #1
}

\newtcolorbox{warningbox}[1][]{
    colback=warningbox,
    colframe=warningred,
    fonttitle=\bfseries,
    title=\faExclamationTriangle\ Warning,
    breakable,
    enhanced,
    attach boxed title to top left={yshift=-2mm, xshift=5mm},
    boxed title style={colback=warningred, colframe=warningred},
    #1
}

\newtcolorbox{successbox}[1][]{
    colback=green!5,
    colframe=successgreen,
    fonttitle=\bfseries,
    title=\faCheckCircle\ Success,
    breakable,
    enhanced,
    attach boxed title to top left={yshift=-2mm, xshift=5mm},
    boxed title style={colback=successgreen, colframe=successgreen},
    #1
}

\newtcolorbox{commandbox}[1][]{
    colback=codebg,
    colframe=codeframe,
    fonttitle=\bfseries\ttfamily,
    title=\faTerminal\ Command,
    breakable,
    enhanced,
    #1
}

% ============================================
% CODE LISTINGS
% ============================================
\lstdefinestyle{pythonstyle}{
    language=Python,
    basicstyle=\ttfamily\small,
    backgroundcolor=\color{codebg},
    frame=single,
    rulecolor=\color{codeframe},
    keywordstyle=\color{blue}\bfseries,
    commentstyle=\color{gray}\itshape,
    stringstyle=\color{red},
    numbers=left,
    numberstyle=\tiny\color{gray},
    stepnumber=1,
    numbersep=10pt,
    tabsize=4,
    showspaces=false,
    showstringspaces=false,
    breaklines=true,
    breakatwhitespace=true,
    captionpos=b
}

\lstdefinestyle{vbastyle}{
    language=VBScript,
    basicstyle=\ttfamily\small,
    backgroundcolor=\color{codebg},
    frame=single,
    rulecolor=\color{codeframe},
    keywordstyle=\color{blue}\bfseries,
    commentstyle=\color{green}\itshape,
    stringstyle=\color{red},
    numbers=left,
    numberstyle=\tiny\color{gray},
    stepnumber=1,
    numbersep=10pt,
    tabsize=4,
    showspaces=false,
    showstringspaces=false,
    breaklines=true,
    breakatwhitespace=true,
    captionpos=b
}

\lstdefinestyle{shellstyle}{
    basicstyle=\ttfamily\small,
    backgroundcolor=\color{codebg},
    frame=single,
    rulecolor=\color{codeframe},
    keywordstyle=\color{blue}\bfseries,
    commentstyle=\color{gray}\itshape,
    numbers=none,
    showspaces=false,
    showstringspaces=false,
    breaklines=true,
    breakatwhitespace=true
}

% ============================================
% DOCUMENT
% ============================================
\begin{document}

% ============================================
% TITLE PAGE
% ============================================
\begin{titlepage}
    \centering
    \vspace*{2cm}
    
    {\Huge\bfseries\color{primaryblue} Payroll Analytics System\par}
    \vspace{0.5cm}
    {\LARGE User Manual\par}
    \vspace{2cm}
    
    {\Large\itshape Installation, Configuration \& Usage Guide\par}
    \vspace{3cm}
    
    \begin{tcolorbox}[colback=infobox, colframe=primaryblue, width=0.8\textwidth]
        \centering
        \large
        \textbf{Complete Setup Instructions}\\[0.5cm]
        From installation to advanced forecasting\\[0.2cm]
        Includes troubleshooting and best practices
    \end{tcolorbox}
    
    \vfill
    
    {\large MARS Project\\
    Integrated Payroll Management System\\[0.3cm]
    \today\par}
\end{titlepage}

% ============================================
% TABLE OF CONTENTS
% ============================================
\tableofcontents
\newpage

% ============================================
% CHAPTER 1: INTRODUCTION
% ============================================
\section{Introduction}

\subsection{About This Manual}

This user manual provides comprehensive instructions for installing, configuring, and using the Payroll Analytics System. Whether you are a first-time user or an experienced analyst, this guide will help you:

\begin{itemize}[leftmargin=*]
    \item \textbf{Install} all required software dependencies
    \item \textbf{Configure} the system for your environment
    \item \textbf{Process} payroll data from Excel files
    \item \textbf{Generate} forecasts using advanced machine learning models
    \item \textbf{Visualize} insights through interactive dashboards
\end{itemize}

\subsection{System Overview}

The Payroll Analytics System is a Python-based application designed to:

\begin{enumerate}[leftmargin=*]
    \item \textbf{Import \& Map:} Load payroll data from various Excel formats and standardize column names
    \item \textbf{Process \& Transform:} Clean, normalize, and prepare data for analysis
    \item \textbf{Analyze \& Forecast:} Apply 8 predictive models including AR(2)+Ridge regularization
    \item \textbf{Visualize \& Report:} Create 10+ interactive visualizations and export results
\end{enumerate}

\begin{infobox}[title=\faLightbulb\ Key Features]
\begin{itemize}[leftmargin=*]
    \item \textbf{AR(2)+Ridge Forecasting:} Best-in-class model with 4.72\% NRMSE
    \item \textbf{Employee Clustering:} Automatic grouping based on salary patterns
    \item \textbf{Interactive Dashboard:} 10 visualization options with filtering
    \item \textbf{Excel Integration:} VBA control panel for seamless workflow
    \item \textbf{Recursive Forecasting:} 6-month predictions with confidence intervals
\end{itemize}
\end{infobox}

\subsection{Prerequisites}

Before installation, ensure you have:

\begin{itemize}[leftmargin=*]
    \item \textbf{Operating System:} Windows 10/11 (64-bit)
    \item \textbf{Microsoft Excel:} 2016 or later with macro support
    \item \textbf{Administrator Rights:} For installing Python and dependencies
    \item \textbf{Disk Space:} At least 2 GB free space
    \item \textbf{Internet Connection:} For downloading packages
\end{itemize}

\newpage

% ============================================
% CHAPTER 2: INSTALLATION
% ============================================
\section{Installation}

\subsection{Step 1: Download the Project}

\subsubsection{Option A: Clone from Repository (Recommended)}

If you have Git installed:

\begin{commandbox}[title=Git Clone Command]
\begin{lstlisting}[style=shellstyle]
cd C:\Users\YourName\Documents
git clone <repository-url> "Integrated Project"
cd "Integrated Project"
\end{lstlisting}
\end{commandbox}

\subsubsection{Option B: Download ZIP Archive}

\begin{enumerate}[leftmargin=*]
    \item Download the project ZIP file from the repository
    \item Extract to: \texttt{C:\textbackslash Users\textbackslash YourName\textbackslash Documents\textbackslash Integrated Project}
    \item Ensure the folder structure matches:
\end{enumerate}

\begin{lstlisting}[style=shellstyle, frame=single]
Integrated Project/
├── data/
│   └── raw/
├── docs/
├── outputs/
├── scripts/
├── src/
│   ├── core/
│   ├── excel_tools/
│   └── features/
├── tests/
└── Control_Panel_New.xlsm
\end{lstlisting}

\begin{warningbox}[title=Important: Folder Location]
\textbf{Do not use paths with special characters or spaces in parent folders.}\\[0.3cm]
\faCheck\ Good: \texttt{C:\textbackslash Users\textbackslash John\textbackslash Documents\textbackslash Integrated Project}\\
\faTimes\ Bad: \texttt{C:\textbackslash Users\textbackslash John Smith\textbackslash My Files\textbackslash Project}
\end{warningbox}

\subsection{Step 2: Install Python 3.10}

\subsubsection{Download Python}

\begin{enumerate}[leftmargin=*]
    \item Visit: \url{https://www.python.org/downloads/}
    \item Download \textbf{Python 3.10.x} (64-bit recommended)
    \item Run the installer
\end{enumerate}

\subsubsection{Installation Options}

\begin{warningbox}[title=Critical: Installation Settings]
During installation, you \textbf{MUST}:
\begin{itemize}
    \item \faCheckSquare\ Check \textbf{"Add Python 3.10 to PATH"}
    \item \faCheckSquare\ Select \textbf{"Install for all users"} (if available)
    \item \faCheckSquare\ Choose \textbf{"Customize installation"}
    \item \faCheckSquare\ Enable \textbf{"pip"} and \textbf{"tcl/tk and IDLE"}
\end{itemize}
\end{warningbox}

\subsubsection{Verify Installation}

Open PowerShell or Command Prompt and run:

\begin{commandbox}[title=Verification Commands]
\begin{lstlisting}[style=shellstyle]
python --version
# Expected output: Python 3.10.x

python -c "import sys; print(sys.executable)"
# Expected output: C:\Users\YourName\AppData\Local\Programs\Python\Python310\python.exe
\end{lstlisting}
\end{commandbox}

\begin{successbox}[title=Success Indicator]
If both commands execute without errors, Python is correctly installed!
\end{successbox}

\newpage

\subsection{Step 3: Install Required Python Packages}

\subsubsection{Navigate to Project Directory}

\begin{commandbox}[title=Change Directory]
\begin{lstlisting}[style=shellstyle]
cd "C:\Users\YourName\Documents\Integrated Project"
\end{lstlisting}
\end{commandbox}

\subsubsection{Install Core Dependencies}

Run the following command to install all required packages:

\begin{commandbox}[title=Package Installation]
\begin{lstlisting}[style=shellstyle]
python -m pip install --upgrade pip
pip install pandas numpy matplotlib scikit-learn openpyxl xlwings
\end{lstlisting}
\end{commandbox}

\subsubsection{Package Descriptions}

\begin{table}[h]
\centering
\begin{tabular}{@{} l p{8cm} @{}}
\toprule
\textbf{Package} & \textbf{Purpose} \\
\midrule
\texttt{pandas} & Data manipulation and analysis \\
\texttt{numpy} & Numerical computations and arrays \\
\texttt{matplotlib} & Data visualization and plotting \\
\texttt{scikit-learn} & Machine learning models (Ridge, Lasso) \\
\texttt{openpyxl} & Reading/writing Excel files (.xlsx) \\
\texttt{xlwings} & Excel-Python integration (VBA bridge) \\
\bottomrule
\end{tabular}
\caption{Required Python packages and their functions}
\end{table}

\begin{infobox}[title=Installation Time]
Package installation typically takes 2-5 minutes depending on your internet connection.
\end{infobox}

\subsubsection{Verify Package Installation}

\begin{commandbox}[title=Verification Script]
\begin{lstlisting}[style=shellstyle]
python -c "import pandas, numpy, matplotlib, sklearn, openpyxl, xlwings; print('All packages installed successfully!')"
\end{lstlisting}
\end{commandbox}

If you see the success message, proceed to configuration.

\begin{warningbox}[title=Troubleshooting: Import Errors]
If you encounter \texttt{ModuleNotFoundError}:
\begin{enumerate}
    \item Ensure you're using the correct Python executable
    \item Try reinstalling the specific package: \texttt{pip install --force-reinstall <package-name>}
    \item Check for typos in package names
\end{enumerate}
\end{warningbox}

\newpage

% ============================================
% CHAPTER 3: CONFIGURATION
% ============================================
\section{Configuration}

This section guides you through configuring all \texttt{FIX\_ME} paths in the codebase to match your local environment.

\begin{warningbox}[title=Critical Step]
\textbf{All configurations in this section are MANDATORY.}\\
The application will not work without updating these paths.
\end{warningbox}

\subsection{Configuration Overview}

You need to update \textbf{11 configuration points} across 8 files:

\begin{table}[h]
\centering
\small
\begin{tabular}{@{} l c l @{}}
\toprule
\textbf{File} & \textbf{Count} & \textbf{Type} \\
\midrule
Control\_Panel\_VBA\_Module.bas & 2 & VBA Constants \\
src/core/config.py & 2 & Python Paths \\
test\_monthly\_mapping\_engine.py & 2 & Absolute Paths \\
dashboard\_gui.py & 1 & Data Path \\
run\_load\_data.py & 1 & Data Path \\
run\_mapped\_data\_loader.py & 1 & Data Path \\
train\_test\_only\_evaluation.py & 1 & Data Path \\
simplified\_forecasting\_comparison.py & 1 & Data Path \\
\bottomrule
\end{tabular}
\caption{Configuration files requiring updates}
\end{table}

\subsection{VBA Configuration (Control Panel)}

\subsubsection{Find Your Python Path}

First, determine your Python executable location:

\begin{commandbox}[title=Get Python Path]
\begin{lstlisting}[style=shellstyle]
python -c "import sys; print(sys.executable)"
\end{lstlisting}
\end{commandbox}

\textbf{Example output:}
\begin{lstlisting}[style=shellstyle, frame=none]
C:\Users\Morta\AppData\Local\Programs\Python\Python310\python.exe
\end{lstlisting}

Copy this entire path for the next step.

\subsubsection{Edit VBA Module}

\begin{enumerate}[leftmargin=*]
    \item Open \texttt{Control\_Panel\_New.xlsm} in Excel
    \item Press \texttt{Alt+F11} to open the VBA Editor
    \item In the Project Explorer, locate \texttt{PayrollControlPanel} module
    \item Find lines 40-41 (marked with \texttt{FIX\_ME} comment)
\end{enumerate}

\begin{lstlisting}[style=vbastyle, caption=VBA Configuration Lines 40-41]
' FIX_ME: Update these paths to match your environment
Const PYTHON_PATH As String = "C:\Users\YourName\...\python.exe"
Const PROJECT_PATH As String = "C:\Users\YourName\...\Integrated Project\"
\end{lstlisting}

\subsubsection{Update the Paths}

Replace with your actual paths:

\begin{lstlisting}[style=vbastyle, caption=Example Configuration]
Const PYTHON_PATH As String = "C:\Users\Morta\AppData\Local\Programs\Python\Python310\python.exe"
Const PROJECT_PATH As String = "C:\Users\Morta\Documents\Integrated Project\"
\end{lstlisting}

\begin{warningbox}[title=Path Format Rules]
\begin{itemize}
    \item Use backslashes (\textbackslash) not forward slashes (/)
    \item \texttt{PROJECT\_PATH} \textbf{must} end with a backslash (\textbackslash)
    \item Use full absolute paths (not relative paths like \texttt{../ or ./})
    \item Enclose paths in double quotes
\end{itemize}
\end{warningbox}

\subsubsection{Test Configuration}

\begin{enumerate}[leftmargin=*]
    \item In the VBA Editor, press \texttt{F5} to run \texttt{VerifyConfiguration} macro
    \item A message box will appear showing:
    \begin{itemize}
        \item \faCheck\ Python found/not found
        \item \faCheck\ Project folder found/not found
    \end{itemize}
    \item If both checks pass, save and close the VBA Editor
\end{enumerate}

\begin{successbox}[title=Configuration Successful]
\textbf{Checkmarks on both items?} Your VBA configuration is correct!
\end{successbox}

\newpage

\subsection{Python Configuration}

\subsubsection{Core Configuration: config.py}

Open \texttt{src/core/config.py} in a text editor.

\textbf{Configuration Point 1: Payroll File Name (Line 18)}

\begin{lstlisting}[style=pythonstyle, caption=Update Payroll File Name]
# FIX_ME: Update the filename to match your actual payroll Excel file
# This should be the main payroll workbook in the data/raw/ directory
PAYROLL_FILE = os.path.join(DATA_RAW_DIR, "Pay emploier sept25.xlsm")
\end{lstlisting}

\textbf{Change to:} Your actual Excel filename in \texttt{data/raw/} folder.

\textbf{Example:}
\begin{lstlisting}[style=pythonstyle]
PAYROLL_FILE = os.path.join(DATA_RAW_DIR, "Control_Panel.xlsm")
# or
PAYROLL_FILE = os.path.join(DATA_RAW_DIR, "company_payroll_2025.xlsx")
\end{lstlisting}

\textbf{Configuration Point 2: Sheet Names (Lines 22-24)}

\begin{lstlisting}[style=pythonstyle, caption=Update Sheet Names]
# FIX_ME: Update sheet names to match your Excel workbook structure
SRC_SHEET = "Feuil1"       # source sheet with full payroll data
TGT_SHEET = "Sheet1"       # simplified / target sheet
CUSTOM_SHEET = "CustomMap" # sheet for saved mappings (auto-created)
\end{lstlisting}

\begin{infobox}[title=How to Find Sheet Names]
\begin{enumerate}
    \item Open your Excel file
    \item Look at the sheet tabs at the bottom
    \item Copy the exact names (case-sensitive)
\end{enumerate}
\end{infobox}

\textbf{Example:} If your sheets are named "Payroll\_Data" and "Summary":
\begin{lstlisting}[style=pythonstyle]
SRC_SHEET = "Payroll_Data"
TGT_SHEET = "Summary"
CUSTOM_SHEET = "CustomMap"  # Keep this as-is (will be created automatically)
\end{lstlisting}

\subsubsection{Mapping Engine: test\_monthly\_mapping\_engine.py}

Open \texttt{tests/test\_monthly\_mapping\_engine.py} in a text editor.

\textbf{Configuration Point: File Paths (Lines 50-51)}

\begin{lstlisting}[style=pythonstyle, caption=Update Absolute Paths]
# FIX_ME: SOURCE AND OUTPUT EXCEL PATHS
SOURCE_EXCEL_PATH = r"C:\Users\Morta\OneDrive\Desktop\Gam3a\MARS\Integrated Project\data\raw\essai donnees paye.xlsx"
OUTPUT_EXCEL_PATH = r"C:\Users\Morta\OneDrive\Desktop\Gam3a\MARS\Integrated Project\data\raw\data_sorted.xlsx"
\end{lstlisting}

\begin{warningbox}[title=Critical: Absolute Paths Required]
These \textbf{MUST} be full absolute paths matching your project location.\\[0.3cm]
\textbf{Use raw strings} (prefix with \texttt{r}) to avoid escape character issues.
\end{warningbox}

\textbf{How to build your paths:}
\begin{enumerate}
    \item Find your project root folder (e.g., \texttt{C:\textbackslash Users\textbackslash John\textbackslash Documents\textbackslash Integrated Project})
    \item Append \texttt{\textbackslash data\textbackslash raw\textbackslash your\_file.xlsx}
    \item Prefix with \texttt{r} to create a raw string
\end{enumerate}

\textbf{Example for user "John" with project in Documents:}
\begin{lstlisting}[style=pythonstyle]
SOURCE_EXCEL_PATH = r"C:\Users\John\Documents\Integrated Project\data\raw\payroll.xlsx"
OUTPUT_EXCEL_PATH = r"C:\Users\John\Documents\Integrated Project\data\raw\data_sorted.xlsx"
\end{lstlisting}

\newpage

\subsubsection{Dashboard: dashboard\_gui.py}

Open \texttt{src/features/analytics/dashboard\_gui.py} in a text editor.

\textbf{Configuration Point: Data Path (Line 40)}

\begin{lstlisting}[style=pythonstyle, caption=Dashboard Data Path]
# FIX_ME: Data file path - Update if your payroll_long.csv is located elsewhere
self.data_path = project_root / "outputs" / "payroll_long.csv"
\end{lstlisting}

\begin{infobox}[title=Default Path]
This uses a \textbf{relative path} that works automatically if your project structure is correct.\\[0.2cm]
\textbf{Only change this if you moved the outputs folder to a different location.}
\end{infobox}

\subsubsection{Data Loading: run\_load\_data.py}

Open \texttt{scripts/run\_load\_data.py} in a text editor.

\textbf{Configuration Point: Workbook Path (Line 14)}

\begin{lstlisting}[style=pythonstyle, caption=Multi-Month Workbook Path]
# FIX_ME: Update this path to point to your multi-month payroll workbook
wb_path = Path("data/raw/synthetic_payroll_startup.xlsx")
\end{lstlisting}

\textbf{Change to:} Your actual multi-month workbook filename.

\textbf{Example:}
\begin{lstlisting}[style=pythonstyle]
wb_path = Path("data/raw/payroll_2025.xlsx")
\end{lstlisting}

\subsubsection{Evaluation Scripts Configuration}

The following scripts use the same configuration pattern:

\begin{itemize}
    \item \texttt{tests/train\_test\_only\_evaluation.py} (Lines 36-38)
    \item \texttt{tests/simplified\_forecasting\_comparison.py} (Lines 47-50)
    \item \texttt{tests/honest\_forecasting\_evaluation.py} (Lines 37-39)
    \item \texttt{tests/run\_mapped\_data\_loader.py} (Lines 29-30)
\end{itemize}

\textbf{Configuration Pattern:}

\begin{lstlisting}[style=pythonstyle, caption=Standard Evaluation Configuration]
# FIX_ME: Configuration - Update data path if needed
TARGET_LABEL = "salaire_brut"  # Column containing salary data
CSV_PATH = "outputs/payroll_long.csv"  # Processed data location
TEST_MONTHS = 6  # Number of months for testing
\end{lstlisting}

\begin{infobox}[title=When to Change These]
These configurations use \textbf{relative paths} and work automatically.\\[0.2cm]
\textbf{Only update if you:}
\begin{itemize}
    \item Renamed the \texttt{outputs} folder
    \item Changed the target column name
    \item Want a different test period length
\end{itemize}
\end{infobox}

\subsection{Configuration Checklist}

Before proceeding, verify you've updated:

\begin{itemize}[label=\faSquare]
    \item[\faCheckSquare] VBA: \texttt{PYTHON\_PATH} constant
    \item[\faCheckSquare] VBA: \texttt{PROJECT\_PATH} constant
    \item[\faCheckSquare] Python: \texttt{PAYROLL\_FILE} in config.py
    \item[\faCheckSquare] Python: Sheet names in config.py
    \item[\faCheckSquare] Python: Absolute paths in test\_monthly\_mapping\_engine.py
    \item[\faSquare] Optional: Dashboard data path (only if moved)
    \item[\faSquare] Optional: Evaluation script paths (only if renamed)
\end{itemize}

\begin{successbox}[title=Configuration Complete]
If you've checked all mandatory items (\faCheckSquare), your system is configured correctly!
\end{successbox}

\newpage

% ============================================
% CHAPTER 4: USING THE SYSTEM
% ============================================
\section{Using the Payroll Analytics System}

\subsection{System Workflow Overview}

The typical workflow consists of 4 main steps:

\begin{enumerate}[leftmargin=*]
    \item \textbf{Import \& Map:} Load raw payroll Excel files and define column mappings
    \item \textbf{Process Data:} Convert to standardized CSV format
    \item \textbf{Analyze:} Run forecasting models and generate predictions
    \item \textbf{Visualize:} Explore insights through the interactive dashboard
\end{enumerate}

\begin{figure}[h]
\centering
\begin{tcolorbox}[width=0.9\textwidth, colback=white, colframe=primaryblue]
\centering
\textbf{System Workflow}\\[0.5cm]
Raw Excel Files $\xrightarrow{\text{Import \& Map}}$ Standardized Data $\xrightarrow{\text{Process}}$ Clean CSV\\[0.3cm]
Clean CSV $\xrightarrow{\text{Analyze}}$ Predictions $\xrightarrow{\text{Visualize}}$ Dashboard
\end{tcolorbox}
\caption{Data processing pipeline}
\end{figure}

\subsection{Opening the Control Panel}

\begin{enumerate}[leftmargin=*]
    \item Navigate to your project folder
    \item Open \texttt{Control\_Panel\_New.xlsm}
    \item If prompted, click \textbf{"Enable Macros"} or \textbf{"Enable Content"}
\end{enumerate}

\begin{warningbox}[title=Macro Security Settings]
If macros are blocked:
\begin{enumerate}
    \item File $\rightarrow$ Options $\rightarrow$ Trust Center $\rightarrow$ Trust Center Settings
    \item Macro Settings $\rightarrow$ Select \textbf{"Enable all macros"}
    \item Restart Excel
\end{enumerate}
\end{warningbox}

The Control Panel provides 4 main buttons:

\begin{table}[h]
\centering
\begin{tabular}{@{} l p{9cm} @{}}
\toprule
\textbf{Button} & \textbf{Function} \\
\midrule
Import \& Map Data & Launch the mapping GUI to process Excel files \\
Open Output Folder & View generated CSV files and results \\
Open Mapped Workbook & Review the standardized Excel output \\
Visualize Data & Launch the analytics dashboard \\
\bottomrule
\end{tabular}
\caption{Control Panel buttons}
\end{table}

\newpage

\subsection{Step 1: Import and Map Data}

\subsubsection{Prepare Your Data}

Before importing, ensure your payroll Excel file:

\begin{itemize}[leftmargin=*]
    \item Is located in \texttt{data/raw/} folder
    \item Contains monthly payroll sheets (e.g., "January", "February", etc.)
    \item Has consistent column headers across all sheets
    \item Includes employee identifiers and salary information
\end{itemize}

\subsubsection{Launch the Mapping GUI}

\begin{enumerate}[leftmargin=*]
    \item In the Control Panel, click \textbf{"Import \& Map Data"}
    \item A Python window will open showing the Import Dialog
    \item The GUI will display all Excel files in your \texttt{data/raw/} folder
\end{enumerate}

\begin{infobox}[title=First-Time Setup]
The first time you run this, it may take 5-10 seconds for Python to initialize.
\end{infobox}

\subsubsection{Select Source File}

\begin{enumerate}[leftmargin=*]
    \item In the file list, select your payroll Excel file
    \item Click \textbf{"Import Selected File"}
    \item The system will analyze the file structure
\end{enumerate}

\subsubsection{Define Column Mappings}

The mapping window shows:

\begin{itemize}[leftmargin=*]
    \item \textbf{Left column:} Target field names (standardized)
    \item \textbf{Right column:} Dropdown menus with your Excel columns
\end{itemize}

\textbf{Required mappings:}

\begin{table}[h]
\centering
\small
\begin{tabular}{@{} l p{7cm} @{}}
\toprule
\textbf{Target Field} & \textbf{Description} \\
\midrule
employee\_id & Unique identifier for each employee \\
salaire\_brut & Gross salary (before deductions) \\
cot\_salarie & Employee social contributions \\
cot\_patronale & Employer social contributions \\
net\_imposable & Taxable net salary \\
PAS & Income tax withheld (Pay-As-You-Earn) \\
net\_paye & Net salary paid to employee \\
avantages & Benefits and allowances \\
total\_cost & Total cost to employer \\
\bottomrule
\end{tabular}
\caption{Standard target fields}
\end{table}

\textbf{Mapping instructions:}

\begin{enumerate}[leftmargin=*]
    \item For each target field, select the corresponding column from your Excel file
    \item If a field doesn't exist in your data, select \textbf{"(none)"}
    \item For calculated fields (e.g., \texttt{total\_cost}), you can use formulas
\end{enumerate}

\begin{infobox}[title=Formula Support]
You can create calculated fields using operators:\\[0.2cm]
\texttt{total\_cost = salaire\_brut + cot\_patronale + avantages}\\[0.2cm]
Click the \textbf{"Formula Builder"} button for assistance.
\end{infobox}

\subsubsection{Save and Apply Mappings}

\begin{enumerate}[leftmargin=*]
    \item After defining all mappings, click \textbf{"Save Mappings"}
    \item The system will:
    \begin{itemize}
        \item Save mappings to a \texttt{CustomMap} sheet (for reuse)
        \item Process all monthly sheets
        \item Create \texttt{data\_sorted.xlsx} with standardized data
        \item Generate \texttt{payroll\_long.csv} for analysis
    \end{itemize}
    \item A success message will appear when processing is complete
\end{enumerate}

\begin{successbox}[title=Processing Complete]
\textbf{Files created:}
\begin{itemize}
    \item \texttt{data/raw/data\_sorted.xlsx} - Standardized Excel workbook
    \item \texttt{outputs/payroll\_long.csv} - Analysis-ready CSV file
\end{itemize}
\end{successbox}

\subsubsection{Verify Output}

\begin{enumerate}[leftmargin=*]
    \item Click \textbf{"Open Output Folder"} in the Control Panel
    \item Check that \texttt{payroll\_long.csv} exists
    \item Open the file to verify:
    \begin{itemize}
        \item All months are included
        \item Column names are standardized
        \item Data values are correct
    \end{itemize}
\end{enumerate}

\newpage

\subsection{Step 2: Review Mapped Data (Optional)}

\subsubsection{Open the Mapped Workbook}

\begin{enumerate}[leftmargin=*]
    \item In the Control Panel, click \textbf{"Open Mapped Workbook"}
    \item Excel will open \texttt{data\_sorted.xlsx}
    \item Review the standardized format
\end{enumerate}

\subsubsection{What to Check}

\begin{itemize}[leftmargin=*]
    \item \textbf{Structure:} Each month is a separate sheet
    \item \textbf{Headers:} Row 1 contains standardized column names
    \item \textbf{Data:} Rows 2+ contain employee records
    \item \textbf{Mappings:} CustomMap sheet shows your saved mappings
\end{itemize}

\begin{infobox}[title=Editing Mappings]
To modify mappings:
\begin{enumerate}
    \item Delete rows from the CustomMap sheet
    \item Re-run "Import \& Map Data"
    \item The GUI will use existing mappings for unchanged fields
\end{enumerate}
\end{infobox}

\subsection{Step 3: Visualize and Analyze}

\subsubsection{Launch the Dashboard}

\begin{enumerate}[leftmargin=*]
    \item In the Control Panel, click \textbf{"Visualize Data"}
    \item The Payroll Analytics Dashboard will open
    \item Wait 2-3 seconds for data to load
\end{enumerate}

\subsubsection{Dashboard Interface}

The dashboard has 3 main sections:

\begin{enumerate}[leftmargin=*]
    \item \textbf{Top Panel:} Filters for employee and month selection
    \item \textbf{Left Panel:} 10 visualization options
    \item \textbf{Right Panel:} Interactive charts and graphs
\end{enumerate}

\subsubsection{Available Visualizations}

\begin{table}[h]
\centering
\small
\begin{tabular}{@{} l p{8cm} @{}}
\toprule
\textbf{Visualization} & \textbf{Description} \\
\midrule
Total Cost Trend & Monthly total payroll cost over time \\
Forecasting & 6-month predictions using AR(2)+Ridge model \\
Cost per Employee & Individual employee cost breakdown \\
Employee Clustering & Automatic grouping based on salary patterns \\
Distribution Analysis & Salary distribution histograms \\
Cost Components & Breakdown of gross, net, and contributions \\
Monthly Comparison & Side-by-side month comparison \\
Correlation Matrix & Relationships between payroll variables \\
Top Earners & Highest-paid employees ranking \\
Cost Variance & Month-to-month cost changes \\
\bottomrule
\end{tabular}
\caption{Dashboard visualizations}
\end{table}

\newpage

\subsubsection{Using Filters}

\textbf{Employee Filter:}
\begin{itemize}[leftmargin=*]
    \item Select \textbf{"All"} to view aggregate data
    \item Select a specific employee to view individual data
\end{itemize}

\textbf{Month Filter:}
\begin{itemize}[leftmargin=*]
    \item Select \textbf{"All"} to include all available months
    \item Select a specific month to focus on that period
\end{itemize}

\begin{infobox}[title=Filter Interactions]
Filters apply to \textbf{all visualizations}. Change filters once and all charts update automatically.
\end{infobox}

\subsubsection{Forecasting Feature}

The \textbf{Forecasting} visualization is the most advanced feature:

\begin{itemize}[leftmargin=*]
    \item Uses \textbf{AR(2)+Ridge} model (best performer: 4.72\% NRMSE)
    \item Predicts 6 months into the future
    \item Shows \textbf{95\% confidence intervals} with error amplification
    \item Includes historical data for context
\end{itemize}

\textbf{How to interpret:}
\begin{enumerate}[leftmargin=*]
    \item \textbf{Blue line:} Historical payroll costs
    \item \textbf{Red dashed line:} Forecasted costs
    \item \textbf{Shaded area:} 95\% confidence interval (uncertainty range)
\end{enumerate}

\textbf{Confidence interval formula:}
\begin{equation}
\text{CI}_t = \pm \text{RMSE} \times \sqrt{t} \times 1.96
\end{equation}

Where $t$ is the forecast step (1-6 months).

\begin{warningbox}[title=Forecast Limitations]
\textbf{Forecasts assume:}
\begin{itemize}
    \item No major structural changes (hiring/layoffs)
    \item Consistent payroll policies
    \item Historical patterns continue
\end{itemize}
Use forecasts as \textbf{guidance}, not absolute predictions.
\end{warningbox}

\subsubsection{Employee Clustering}

The \textbf{Clustering} visualization automatically groups employees:

\begin{itemize}[leftmargin=*]
    \item Groups based on salary patterns and trends
    \item Uses K-Means algorithm with optimal cluster count
    \item Shows cluster centroids (average profiles)
\end{itemize}

\textbf{Use cases:}
\begin{itemize}[leftmargin=*]
    \item Identify salary tiers (junior/mid/senior)
    \item Detect outliers or anomalies
    \item Plan compensation strategies
\end{itemize}

\newpage

\subsection{Step 4: Advanced Analysis (Optional)}

\subsubsection{Running Model Evaluations}

To compare all 8 forecasting models:

\begin{commandbox}[title=Run Model Evaluation]
\begin{lstlisting}[style=shellstyle]
cd "C:\Users\YourName\Documents\Integrated Project"
python tests/train_test_only_evaluation.py
\end{lstlisting}
\end{commandbox}

This will:
\begin{enumerate}[leftmargin=*]
    \item Evaluate all models (AR(1), AR(2), AR(2)+Ridge, AR(2)+Lasso, Ridge-Panel, Lasso-Panel, Pooled FE, Unpooled)
    \item Use Train/Test split and Recursive multi-step methods
    \item Save results to \texttt{outputs/train\_test\_only\_evaluation.csv}
    \item Display RMSE, NRMSE, MAE, and $R^2$ metrics
\end{enumerate}

\begin{infobox}[title=Evaluation Time]
Model evaluation takes approximately 5-10 minutes for 30 employees × 24 months.
\end{infobox}

\subsubsection{Model Performance Summary}

\begin{table}[h]
\centering
\small
\begin{tabular}{@{} l c c c @{}}
\toprule
\textbf{Model} & \textbf{Train/Test NRMSE} & \textbf{Recursive NRMSE} & \textbf{Rank} \\
\midrule
AR(2)+Ridge & 0.0472 & 0.0699 & \textbf{1} \\
AR(2) & 0.0534 & 0.0718 & 2 \\
AR(2)+Lasso & 0.0577 & 0.0731 & 3 \\
AR(1) & 0.0616 & 0.0750 & 4 \\
Unpooled & 0.0836 & 0.0836 & 5 \\
Pooled FE & 0.0857 & 0.0857 & 6 \\
Lasso-Panel & 0.0859 & 0.0862 & 7 \\
Ridge-Panel & 0.0863 & 0.0867 & 8 \\
\bottomrule
\end{tabular}
\caption{Model comparison (lower NRMSE is better)}
\end{table}

\subsubsection{Generating LaTeX Tables}

To create publication-ready tables for reports:

\begin{commandbox}[title=Generate LaTeX Tables]
\begin{lstlisting}[style=shellstyle]
python tests/generate_forecast_error_table.py
\end{lstlisting}
\end{commandbox}

Output files (in \texttt{outputs/}):
\begin{itemize}[leftmargin=*]
    \item \texttt{forecasting\_complete\_table.tex} - Full model comparison
    \item \texttt{forecast\_error\_propagation\_table.tex} - Error amplification
\end{itemize}

These can be directly included in \LaTeX documents.

\newpage

% ============================================
% CHAPTER 5: TROUBLESHOOTING
% ============================================
\section{Troubleshooting}

\subsection{Python Issues}

\subsubsection{Problem: "Python is not recognized"}

\textbf{Symptoms:}
\begin{lstlisting}[style=shellstyle]
'python' is not recognized as an internal or external command
\end{lstlisting}

\textbf{Solutions:}
\begin{enumerate}[leftmargin=*]
    \item Verify Python is installed: Open \texttt{Add/Remove Programs} and search "Python"
    \item Add Python to PATH:
    \begin{enumerate}
        \item Search for "Environment Variables" in Windows
        \item Edit "Path" variable
        \item Add: \texttt{C:\textbackslash Users\textbackslash YourName\textbackslash AppData\textbackslash Local\textbackslash Programs\textbackslash Python\textbackslash Python310}
        \item Add: \texttt{C:\textbackslash Users\textbackslash YourName\textbackslash AppData\textbackslash Local\textbackslash Programs\textbackslash Python\textbackslash Python310\textbackslash Scripts}
        \item Restart Command Prompt/PowerShell
    \end{enumerate}
    \item Use full path: Instead of \texttt{python}, use the full executable path
\end{enumerate}

\subsubsection{Problem: "No module named 'pandas'"}

\textbf{Symptoms:}
\begin{lstlisting}[style=shellstyle]
ModuleNotFoundError: No module named 'pandas'
\end{lstlisting}

\textbf{Solutions:}
\begin{enumerate}[leftmargin=*]
    \item Verify you're using the correct Python:
    \begin{lstlisting}[style=shellstyle]
python -c "import sys; print(sys.executable)"
    \end{lstlisting}
    \item Reinstall the package:
    \begin{lstlisting}[style=shellstyle]
pip install --force-reinstall pandas
    \end{lstlisting}
    \item Check if pip is installing to the correct location:
    \begin{lstlisting}[style=shellstyle]
pip --version
    \end{lstlisting}
\end{enumerate}

\subsection{Excel/VBA Issues}

\subsubsection{Problem: "Macros are disabled"}

\textbf{Symptoms:}
\begin{itemize}[leftmargin=*]
    \item Buttons don't work
    \item Yellow warning banner at top of Excel
\end{itemize}

\textbf{Solutions:}
\begin{enumerate}[leftmargin=*]
    \item Click \textbf{"Enable Content"} in the warning banner
    \item OR: File $\rightarrow$ Options $\rightarrow$ Trust Center $\rightarrow$ Trust Center Settings $\rightarrow$ Macro Settings $\rightarrow$ Enable all macros
\end{enumerate}

\subsubsection{Problem: "Run-time error '53': File not found"}

\textbf{Symptoms:}
\begin{itemize}[leftmargin=*]
    \item VBA error when clicking buttons
    \item Error points to \texttt{PYTHON\_PATH} or \texttt{PROJECT\_PATH}
\end{itemize}

\textbf{Solutions:}
\begin{enumerate}[leftmargin=*]
    \item Verify paths in VBA are correct (see Section 3.2)
    \item Run \texttt{VerifyConfiguration} macro
    \item Check for typos in path strings
    \item Ensure backslashes, not forward slashes
\end{enumerate}

\subsection{Data Processing Issues}

\subsubsection{Problem: "No data available"}

\textbf{Symptoms:}
\begin{itemize}[leftmargin=*]
    \item Dashboard opens but shows empty charts
    \item Error: "No data available for visualization"
\end{itemize}

\textbf{Solutions:}
\begin{enumerate}[leftmargin=*]
    \item Verify \texttt{outputs/payroll\_long.csv} exists
    \item Check CSV file is not empty (should have multiple rows)
    \item Re-run "Import \& Map Data" from Control Panel
    \item Verify mappings were saved correctly
\end{enumerate}

\subsubsection{Problem: "Insufficient data for forecasting"}

\textbf{Symptoms:}
\begin{itemize}[leftmargin=*]
    \item Error: "Need at least 5 months of data"
\end{itemize}

\textbf{Solutions:}
\begin{enumerate}[leftmargin=*]
    \item Ensure your Excel file contains at least 5 monthly sheets
    \item Check that all months have data (not empty sheets)
    \item Verify employee IDs are consistent across months
\end{enumerate}

\subsection{Performance Issues}

\subsubsection{Problem: "Application is slow"}

\textbf{Symptoms:}
\begin{itemize}[leftmargin=*]
    \item Dashboard takes long to open
    \item Visualizations lag when changing filters
\end{itemize}

\textbf{Solutions:}
\begin{enumerate}[leftmargin=*]
    \item \textbf{Reduce data size:} Filter to specific date ranges
    \item \textbf{Close other applications:} Free up RAM
    \item \textbf{Use SSD:} Move project to solid-state drive
    \item \textbf{Upgrade Python packages:} \texttt{pip install --upgrade pandas numpy}
\end{enumerate}

\subsection{Getting Help}

If you encounter issues not covered here:

\begin{enumerate}[leftmargin=*]
    \item Check the \texttt{README.md} file in the project root
    \item Review error messages carefully - they often indicate the problem
    \item Check Python/VBA syntax for typos in configuration files
    \item Consult the project documentation in \texttt{docs/} folder
\end{enumerate}

\newpage

% ============================================
% CHAPTER 6: BEST PRACTICES
% ============================================
\section{Best Practices}

\subsection{Data Management}

\subsubsection{Organize Your Files}

\begin{itemize}[leftmargin=*]
    \item \textbf{Keep raw data separate:} Place original Excel files in \texttt{data/raw/}
    \item \textbf{Version your data:} Use date suffixes (e.g., \texttt{payroll\_2025\_01.xlsx})
    \item \textbf{Backup regularly:} Copy \texttt{outputs/} folder to external storage
    \item \textbf{Document changes:} Maintain a log of data processing steps
\end{itemize}

\subsubsection{Data Quality Checks}

Before importing:
\begin{enumerate}[leftmargin=*]
    \item \textbf{Verify completeness:} All required columns present
    \item \textbf{Check formats:} Numbers stored as numbers, not text
    \item \textbf{Remove duplicates:} No duplicate employee records per month
    \item \textbf{Validate ranges:} Salaries within expected bounds
\end{enumerate}

\subsection{Forecasting Guidelines}

\subsubsection{When to Use Forecasts}

Forecasts are most reliable when:
\begin{itemize}[leftmargin=*]
    \item You have \textbf{at least 12 months} of historical data
    \item Payroll patterns are \textbf{relatively stable}
    \item No recent major organizational changes
    \item Employee turnover is low to moderate
\end{itemize}

\subsubsection{Interpreting Confidence Intervals}

\begin{itemize}[leftmargin=*]
    \item \textbf{Narrow intervals:} High confidence, stable predictions
    \item \textbf{Wide intervals:} High uncertainty, use caution
    \item \textbf{Growing intervals:} Normal - uncertainty increases over time
    \item \textbf{95\% meaning:} 95\% of actual values will fall within the shaded area
\end{itemize}

\subsubsection{Model Selection}

\begin{itemize}[leftmargin=*]
    \item \textbf{Default choice:} AR(2)+Ridge (best overall performance)
    \item \textbf{For stability:} Use recursive forecasting method
    \item \textbf{For accuracy:} Retrain models every 3-6 months
    \item \textbf{For interpretation:} Review NRMSE (normalized metric)
\end{itemize}

\subsection{Visualization Tips}

\subsubsection{Effective Filtering}

\begin{itemize}[leftmargin=*]
    \item Start with \textbf{"All"} filters to see overall trends
    \item Focus on \textbf{specific employees} for detailed analysis
    \item Use \textbf{monthly filters} to investigate anomalies
    \item Compare \textbf{before/after} periods when changes occur
\end{itemize}

\subsubsection{Exporting Results}

To save visualizations:
\begin{enumerate}[leftmargin=*]
    \item Generate the desired chart in the dashboard
    \item Right-click on the chart
    \item Select "Save Image As..."
    \item Choose PNG format for presentations, SVG for publications
\end{enumerate}

\subsection{Maintenance}

\subsubsection{Regular Updates}

\textbf{Monthly tasks:}
\begin{itemize}[leftmargin=*]
    \item Import new payroll data
    \item Review forecast accuracy
    \item Update visualizations
\end{itemize}

\textbf{Quarterly tasks:}
\begin{itemize}[leftmargin=*]
    \item Re-run model evaluation
    \item Archive old data
    \item Backup project folder
\end{itemize}

\textbf{Annual tasks:}
\begin{itemize}[leftmargin=*]
    \item Update Python packages: \texttt{pip install --upgrade -r requirements.txt}
    \item Review and clean up \texttt{outputs/} folder
    \item Document any custom modifications
\end{itemize}

\subsubsection{Version Control}

If using Git:
\begin{lstlisting}[style=shellstyle]
# Save your configurations
git add src/core/config.py
git add Control_Panel_VBA_Module.bas
git commit -m "Updated configuration for environment"

# Don't commit sensitive data
# Add to .gitignore:
data/raw/*.xlsx
outputs/*.csv
\end{lstlisting}

\newpage

% ============================================
% APPENDIX
% ============================================
\section{Appendix}

\subsection{Appendix A: File Structure Reference}

\begin{lstlisting}[style=shellstyle, caption=Complete Project Structure]
Integrated Project/
├── Control_Panel_New.xlsm          # Main Excel control panel
├── Control_Panel_VBA_Module.bas    # VBA code for buttons
├── data/
│   ├── raw/                        # Raw Excel payroll files
│   │   ├── Pay emploier sept25.xlsm
│   │   ├── data_sorted.xlsx        # Standardized output
│   │   └── essai donnees paye.xlsx
│   └── processed/                  # (Optional) Intermediate files
├── docs/                           # Documentation
│   ├── USER_MANUAL.tex             # This manual (LaTeX source)
│   ├── CLUSTERING_GUIDE.md
│   └── latex/                      # Report LaTeX files
├── outputs/                        # Generated results
│   ├── payroll_long.csv            # Main analysis dataset
│   ├── train_test_only_evaluation.csv  # Model results
│   ├── forecasting_complete_table.tex
│   └── forecast_error_propagation_table.tex
├── scripts/                        # Utility scripts
│   ├── run_dashboard.py            # Launch dashboard
│   ├── run_load_data.py            # Load multi-month data
│   └── run_analysis.py
├── src/                            # Source code
│   ├── core/                       # Core utilities
│   │   ├── config.py               # Configuration (FIX_ME)
│   │   ├── data_loader.py
│   │   └── normalization.py
│   ├── excel_tools/                # Excel integration
│   │   ├── import_dialog.py       # Import GUI
│   │   ├── mapping_utils.py
│   │   └── mapping_gui.py
│   └── features/                   # Features
│       └── analytics/
│           ├── dashboard_gui.py    # Main dashboard (FIX_ME)
│           └── employee_clustering.py
└── tests/                          # Analysis scripts
    ├── train_test_only_evaluation.py       # Model evaluation (FIX_ME)
    ├── test_monthly_mapping_engine.py     # Mapping engine (FIX_ME)
    ├── run_mapped_data_loader.py          # Data loader (FIX_ME)
    ├── simplified_forecasting_comparison.py
    └── honest_forecasting_evaluation.py
\end{lstlisting}

\subsection{Appendix B: Configuration Quick Reference}

\begin{table}[h]
\centering
\small
\begin{tabular}{@{} p{5cm} p{8cm} @{}}
\toprule
\textbf{File} & \textbf{What to Change} \\
\midrule
Control\_Panel\_VBA\_Module.bas & Lines 40-41: \texttt{PYTHON\_PATH}, \texttt{PROJECT\_PATH} \\
src/core/config.py & Line 18: \texttt{PAYROLL\_FILE}, Lines 22-24: Sheet names \\
test\_monthly\_mapping\_engine.py & Lines 50-51: Absolute file paths \\
dashboard\_gui.py & Line 40: Data path (optional) \\
run\_load\_data.py & Line 14: Workbook path (optional) \\
\bottomrule
\end{tabular}
\caption{FIX\_ME configuration summary}
\end{table}

\newpage

\subsection{Appendix C: Model Technical Details}

\subsubsection{AR(2)+Ridge Model}

\textbf{Formula:}
\begin{equation}
y_t = \beta_0 + \beta_1 t + \beta_2 \text{month}_t + \beta_3 y_{t-1} + \beta_4 y_{t-2} + \beta_5 y_{t-3} + \epsilon_t
\end{equation}

With L2 regularization:
\begin{equation}
\min_{\beta} \sum_{i=1}^{n} (y_i - \hat{y}_i)^2 + \alpha \sum_{j=1}^{p} \beta_j^2
\end{equation}

\textbf{Parameters:}
\begin{itemize}[leftmargin=*]
    \item $\alpha = 10.0$ (optimal via cross-validation)
    \item Features: time index, month number, 3 lagged values
    \item Training: Per-employee models with standardized features
\end{itemize}

\textbf{Performance:}
\begin{itemize}[leftmargin=*]
    \item RMSE: €128.32 (Train/Test), €190.09 (Recursive)
    \item NRMSE: 0.0472 (Train/Test), 0.0699 (Recursive)
    \item $R^2$: 0.9780 (Train/Test), 0.9517 (Recursive)
\end{itemize}

\subsubsection{Error Propagation}

Recursive forecasting accumulates errors over time:

\begin{equation}
\text{CI}_t = \pm \text{RMSE} \times \sqrt{t} \times 1.96
\end{equation}

\begin{table}[h]
\centering
\begin{tabular}{@{} c c c c @{}}
\toprule
\textbf{Step} & \textbf{Horizon} & \textbf{Error Margin} & \textbf{Growth} \\
\midrule
1 & 1 month & ±€251.53 & -- \\
2 & 2 months & ±€355.72 & +41.4\% \\
3 & 3 months & ±€435.71 & +73.2\% \\
4 & 4 months & ±€503.07 & +100.0\% \\
5 & 5 months & ±€562.35 & +123.6\% \\
6 & 6 months & ±€615.83 & +144.9\% \\
\bottomrule
\end{tabular}
\caption{Forecast error propagation (AR(2)+Ridge, RMSE=€128.32)}
\end{table}

\subsection{Appendix D: Python Package Versions}

Recommended versions (tested and verified):

\begin{table}[h]
\centering
\begin{tabular}{@{} l c @{}}
\toprule
\textbf{Package} & \textbf{Version} \\
\midrule
Python & 3.10.x \\
pandas & $\geq$ 1.5.0 \\
numpy & $\geq$ 1.23.0 \\
matplotlib & $\geq$ 3.6.0 \\
scikit-learn & $\geq$ 1.2.0 \\
openpyxl & $\geq$ 3.0.0 \\
xlwings & $\geq$ 0.28.0 \\
\bottomrule
\end{tabular}
\caption{Minimum package versions}
\end{table}

To check your installed versions:
\begin{lstlisting}[style=shellstyle]
pip list | findstr "pandas numpy matplotlib scikit-learn openpyxl xlwings"
\end{lstlisting}

\newpage

\subsection{Appendix E: Keyboard Shortcuts}

\subsubsection{Excel Shortcuts}

\begin{table}[h]
\centering
\begin{tabular}{@{} l l @{}}
\toprule
\textbf{Shortcut} & \textbf{Action} \\
\midrule
Alt+F11 & Open VBA Editor \\
F5 & Run macro / VBA code \\
Ctrl+G & Open Immediate Window (VBA) \\
Alt+Q & Close VBA Editor \\
Ctrl+S & Save workbook \\
\bottomrule
\end{tabular}
\end{table}

\subsubsection{Dashboard Shortcuts}

\begin{table}[h]
\centering
\begin{tabular}{@{} l l @{}}
\toprule
\textbf{Action} & \textbf{Method} \\
\midrule
Change visualization & Click button in left panel \\
Update filters & Use dropdown menus at top \\
Close dashboard & Close window or press Alt+F4 \\
\bottomrule
\end{tabular}
\end{table}

\subsection{Appendix F: Glossary}

\begin{description}[leftmargin=3cm, style=nextline]
    \item[AR(2)] AutoRegressive model of order 2 - uses 2 previous values to predict next value
    \item[Ridge Regression] Linear regression with L2 regularization penalty to prevent overfitting
    \item[NRMSE] Normalized Root Mean Squared Error - scale-independent accuracy metric
    \item[RMSE] Root Mean Squared Error - average prediction error in original units (€)
    \item[Recursive Forecasting] Multi-step prediction where each forecast feeds into the next
    \item[Confidence Interval] Range expected to contain true value with specified probability (95\%)
    \item[Clustering] Grouping similar employees based on salary patterns
    \item[VBA] Visual Basic for Applications - Excel's macro programming language
    \item[Payroll Long Format] Data structure with one row per employee-month observation
\end{description}

\newpage

% ============================================
% END MATTER
% ============================================

\section*{Document Information}

\begin{table}[h]
\centering
\begin{tabular}{@{} l l @{}}
\toprule
\textbf{Field} & \textbf{Value} \\
\midrule
Title & Payroll Analytics System - User Manual \\
Version & 1.0 \\
Date & \today \\
Project & MARS Integrated Payroll Management \\
Target Audience & System administrators, data analysts, HR managers \\
Document Type & Installation and usage guide \\
\bottomrule
\end{tabular}
\end{table}

\vspace{1cm}

\begin{tcolorbox}[colback=infobox, colframe=primaryblue]
\textbf{Need Help?}\\[0.3cm]
If you encounter issues not covered in this manual:
\begin{itemize}
    \item Review the troubleshooting section (Chapter 5)
    \item Check configuration files for \texttt{FIX\_ME} tags
    \item Verify all paths are correct
    \item Consult project documentation in \texttt{docs/} folder
\end{itemize}
\end{tcolorbox}

\vfill

\begin{center}
\textit{This manual covers the complete setup and usage of the Payroll Analytics System.\\
For technical details about forecasting methodology, see the thesis documentation.}
\end{center}

\end{document}
